\section{One Attempt, and the Patterns Inside It}\label{sec:catalog}

The patterns in this section are the recurring design moves that make up an
active-inference control layer for a coding agent.  Rather than list them, we
walk through a single real run of the system and meet each pattern where it
actually bites.

The run is \emph{attempt-061}.  It begins at 10:46:43 and closes
at 10:59:53 --- thirteen minutes and ten seconds, three agent turns --- and it
ends in the outcome the system is built to produce: a reviewed code change,
written into an external store, that moved a recorded dial.  Every timestamp
below is taken from the run's own durable phase log, not reconstructed for the
telling.  The same attempt is tabulated step-by-step in the companion paper; a
reader who wants the machine-readable version will find it there.

We chose a run that went well, and it is worth saying what that does and does
not buy.  A successful specimen shows what the machinery is \emph{for}.  It does
not show that the machinery usually succeeds --- it usually does not, and the
failure ledger is in the companion paper --- and it does not show that the
system helps.  What it does show, which is the claim this paper makes, is that
each \emph{run-evidenced} move below ran at a stated time and left a record
someone else can check.  The tested entries have concrete code and tests but
are not credited with equivalent run evidence; staged entries are proposals
forced by the record, not presented as completed machinery.\worknote{
attempt-061 remains real and pinned; its cohort
(\texttt{:wm-outer-loop-46-v1}, target 3, closed July) is exactly why the
2026-09-13 machinery run r1 refused admission --- the stopping rule held
rather than minting a fourth attempt. That refusal is this section's
discipline working. A fresh reviewed cohort is staged so new specimens can
exist.}

Each pattern is stated in the \textbf{If} / \textbf{However} / \textbf{Then} /
\textbf{Because} form: a context (\emph{if}), the tension that recurs
(\emph{however}), the move that resolves or ameliorates it (\emph{then}), and
its rationale (\emph{because}).  Most also carry an \emph{invariant to preserve}
--- the one property a different implementation would have to keep for the
pattern to have been applied at all, as opposed to merely cited.  The
parenthetical R-number keys each pattern to the control map in
Figure~\ref{fig:aif-brain}.

Patterns marked \apparatusmark{} are \emph{apparatus}: engineering discipline
that makes the loop trustworthy, rather than active-inference constructs.  They
are kept in the story where they belong rather than quarantined at the end,
because in practice they are what the inference rests on --- but the
distinction is real and worth marking.

\paragraph*{A note on the R-numbers.}
\ifdefined\plopearlyoverview
The R-numbers in Figure~\ref{fig:aif-brain} name the catalogue patterns
directly: an R-number in this paper is a direct join between the figure, the
pattern map, and the pattern text.  The source audit retains the concordance
with an older \emph{completeness contract}, in which R6, R9, and R12 had
narrower or different meanings, but the reader does not have to perform it.  Prior
preferences are the model's C vector rather than an R-numbered pattern, and
R20 is included as the catalogue's interoceptive apparatus.
\else
These are \emph{patterns} --- design moves one might make.  Their numbers were inherited from an existing
\emph{completeness contract} in which the same prefix means something else: a
\emph{requirement}, or a property an implementation must satisfy to carry the
active-inference label honestly.  That contract predates this catalogue, was
written for another project by diffing an implementation against the canonical
literature, and is used unchanged by two other surfaces in our stack.  The
figure retains the contract's node labels; the R-number is therefore a join key,
not a claim that every node and pattern still has the same name.

Most of the original numbered spine remains the same item seen from two sides,
but three labels have drifted and the reader should know which.  At R6 the
contract means softmax selection with abstain, where the catalogue means the
candidate action space; at R12 it means an outer loop inferring the inner loop's
hyperparameters, where the catalogue means two layers of calibration; and at R9
it means named validation properties, where the catalogue means no
self-certification.  The last is the one to watch, because our R9 is the pattern
we would keep if we could keep only one, and because a citation of the contract's
R9 artifact sat in this paper for some months looking entirely correct.  Later
extensions at the learning frontier and R20 sit outside that original
one-to-one concordance.  We report the divergence rather than quietly
renumbering, since the original numbers remain load-bearing in a predecessor
paper; the full concordance is kept alongside the source.
\fi

The catalogue is stated here without its mathematics.  Every quantity named
below --- the belief state, the trust weights, the risk and ambiguity legs of
the action score, the evidence threshold for merging concepts --- has an exact
formal definition, a worked numerical example, and a pointer into the reference
implementation; all of that lives in the companion paper, and nothing in this
section depends on having read it.  The claim here is that the moves are
recognisable and reusable as \emph{design} moves, whatever formalism a later
implementer chooses to express them in.

\ifdefined\plopearlyoverview
\subsection*{Pattern map}

The chronological account deliberately introduces patterns in the order in
which the attempt needed them, not in numerical or conceptual order.  The map
below is a reader's index to that account, not a second catalogue.  A dagger
marks apparatus, and ``staged'' retains the status stated by the corresponding
entry.

\begingroup
\small
\setlength{\tabcolsep}{4pt}
\begin{longtable}{@{}>{\raggedright\arraybackslash}p{.16\linewidth}
                       >{\raggedright\arraybackslash}p{.39\linewidth}
                       >{\raggedright\arraybackslash}p{.39\linewidth}@{}}
\toprule
Category & Pattern ID and name & Brief role \\
\midrule
\endhead
Observation and belief & R1 --- Belief State as Operational Hypotheses & Makes the state updateable and falsifiable. \\
 & R2 --- Structured Observation Vector & Normalises heterogeneous evidence. \\
 & R3 --- Predictive-Coding Belief Update & Moves beliefs by trusted prediction error. \\
 & R4 --- Shared Kernel for the Predictive Forward Model & Keeps predicted and live dynamics aligned. \\
 & R7 --- Evidence Precision Registry & Makes trust explicit per evidence channel. \\
 & R8 --- Present-Fit Mismatch as a Per-Tick Scalar & Exposes how well present belief fits observation. \\
\addlinespace
Evaluation, selection, and planning & R5 --- Expected Free Energy Inside an Auditable Controller & Separates the canonical score from engineering controls. \\
 & R6 --- Candidate Pattern Action Space & Constructs a bounded, reason-bearing choice set. \\
 & R11\apparatusmark{} --- Hierarchical, Budget-Aware Action Selection & Enforces shared budgets at each hierarchy level. \\
 & R13 --- Temporal Depth Beyond the Greedy Step & Scores cascades rather than only their first move. \\
 & R14 --- Selection Gain as Commitment Temperature & Makes exploration versus commitment an explicit dial. \\
 & R15 --- Hierarchical and Temporal Depth & Couples strategic and tactical loops without conflating them. \\
\addlinespace
Actuation and assurance & R9\apparatusmark{} --- No Self-Certification & Requires evidence its claimant did not manufacture. \\
 & R10\apparatusmark{} --- Scheduled Observer Entrypoint & Counts state change, not scheduler firing, as liveness. \\
 & R12 --- Two-Layer Calibration & Separates internal consistency from external value evidence. \\
 & R16 --- Grounded Actuation, Not Re-Observation & Closes the loop with an external typed witness. \\
 & R16 realised\apparatusmark{} --- Semilattice-Rollout Execution Witness & Independently checks that a cascade was constructed. \\
 & R20 --- Interoceptive Tripwires & Makes failures of the harness itself observable. \\
\addlinespace
Learning and structural change & R17 --- Structure Learning by Bayesian Model Reduction & Merges or prunes structure only when evidence warrants it. \\
 & R17$'$ --- Posterior Spread as an Exploration Control & Keeps a useful heuristic honestly outside canonical EFE. \\
 & R17$''$ --- Embedding-Nearness as Generator & Proposes, but does not approve, candidate structural changes. \\
 & R17$'''$ --- Pattern Genesis from Evidence-Bearing Holes & Enlarges the policy vocabulary under witnessed review. \\
 & R16 $\rightarrow$ R6 --- Discharge-Trained Cascade Proposal & Learns diverse proposals from realised outcomes. \\
 & Capability Zones (staged) & Grounds preferences among as-yet-unenacted capabilities. \\
 & Native-Currency Discrimination (staged) & Makes the canonical score discriminate without a bolt-on term. \\
\bottomrule
\end{longtable}
\endgroup
\fi

\subsection*{10:46:43 --- The click, and what kind of click it was}

The attempt begins.  The loop takes a single-flight lock so that two runs cannot
race, and the attempt is accepted.

It matters a great deal \emph{what started it}, and the trace says: this one was
an on-demand click, not a scheduled tick.  The record carries the trigger
explicitly (\texttt{:trigger :duree-click-on-demand}), and it must, because the
scheduled entrypoint the next pattern describes was \emph{operator-disabled}
throughout the events recounted here --- switched off after a stop-line, with
no cron installed.  We flag this at the outset because it would be easy to read
what follows as a portrait of an unattended system running overnight, and it is
not.  It is a portrait of the loop running end to end; a human asked it to
start.

The distinction is one the catalogue itself insists on, and the disabling is the
pattern working rather than failing: the rule below says that when a scheduled
loop cannot change anything, you switch the schedule off rather than let it
write green rows.  That is exactly what the operator did.

\paragraph{Scheduled Observer Entrypoint (R10, the No-Op Trap).\apparatusmark{}} \textbf{If} you want the agent to run live --- to observe, deliberate, and act on a schedule, accumulating trace and precision over real time without a human pressing the button. \textbf{However} a scheduler that fires ticks which cache or no-op (because nothing downstream can change) satisfies the letter of ``live operation'' while the organism is dead, and every firing writes a green ``it ran'' row that hides the inertia. \textbf{Then} install the scheduled entrypoint, but gate ``live'' on evidence that a tick changed observable state (a dial moved, a hole closed), not merely that it fired; when the loop can only no-op, disable the schedule rather than let it paint green. \textbf{Because} liveness is a property of state change, not of the scheduler firing --- a cron over a sealed loop is a clock ticking in an empty room.

\emph{Operational witness:} the hourly entrypoint completed a production-shaped opportunity as a reviewed Git commit plus a Futon1b implementation/discharge, closed the cohort record as \texttt{grounded-change}, and queued an operator QA item.  The scheduler firing alone is not counted as that witness; the external dial movement is.  The hourly schedule was operator-disabled the following day, after a stop-line, and no cron is installed; runs since then --- including the specimen above --- have been on-demand clicks.  We report this rather than the more flattering present tense because the pattern's own rule is that a schedule which cannot change anything should be switched off, and honouring that rule is not a result to be hidden.  The claim this pattern supports is therefore that the loop \emph{can} run unattended end to end and has, not that it is doing so today.\worknote{IMPROVED: the on-demand entrypoint is now a
reviewed, committed client (row 26) rather than an ad-hoc invocation, and
r1 exercised it 2026-09-13. The state-change gate was honoured in the
awkward direction: r1 changed nothing and its record says so
(\texttt{:traceWritten false}, one route hop). \wkflag{OUTSTANDING}: the production
caller that would make R10 credit whole (tracker row 17 names the
integration).}

\subsection*{10:48:48 --- Checking the instruments before trusting them}

Three probes in seven seconds, and the third took thirteen: is the cast of
agents actually invocable; what code is this loop running (revisions and dirty
flags, recorded into provenance before anything else); and does the
authoritative store answer, proved by escalating-patience probes rather than
assumed.  Only when the store answered did the loop claim attempt-061 as its
own.

The ordering is the point.  A system that believes first and checks later has
already contaminated its evidence.  The pattern that governs this phase is the
hardest-won one in the catalogue, and it is the one we would keep if we could
keep only one.

\paragraph{No Self-Certification (R9).\apparatusmark{}} \textbf{If} an instrument issues a verdict from evidence it also produces or stages (a calibrator over its own predictions, a gate over its own outputs). \textbf{However} on first contact with reality every layer tries to certify \emph{itself}, and each attempt looks rigorous --- the forward model scores its own constants (``prior~$=$~value''), the evidence reader eats its own staging files --- and vigilance cannot catch these, because an instrument's first catches are always about itself. \textbf{Then} require that a verdict move \emph{only} on evidence its maker did not manufacture: tag every evidence record at birth with properties the apparatus cannot fake (\texttt{:independent?}, \texttt{:measured} vs \texttt{:target-absent-fallback}, \texttt{:settled} vs \texttt{:transient}) and count only records whose tags clear --- untagged never counts.  The rule covers verdicts about verdicts: a review is
itself a verdict-producing act, and a reviewer's claim that a review ran is
exactly the kind of self-manufactured evidence the pattern bars (measured as
\texttt{:review-execution-evidence-missing}, twice in the first cohort), so
review execution is corroborated from a job-event stream the reviewer does
not control before any review counts. \textbf{Because} circularity is invisible from inside the loop (every internal reading is consistent with every other); only a referent the maker had no hand in breaks it, and structure (``the tag is absent, so it does not count'') survives where vigilance fails.  \emph{Scope in this implementation:} the universal birth-tag regime is the invariant a conforming implementation must keep, not a description of this one.  The running stack has birth tags on the rollout/flight paths and review-execution corroboration conditional on changed code files; it has no catalogue-wide structure that tags every evidence record or excludes every untagged record.  Those narrower mechanisms are evidence of partial application only.\worknote{ACTIVE, applied to itself: the built
no-self-certification checker currently refuses its own admission
(\texttt{:r9/anchor-missing}) until a genesis authority external to it is
authenticated --- the pattern eating its own cooking. Cross-agent genesis
is the chosen anchor design; an external operator-event root has been
located but not yet authenticated. \wkflag{OUTSTANDING}: durable retention of
review commissions (their digest preimages), without which the
review-request join can never pass on real data --- the real-pair run
refused honestly on exactly this.}

\subsection*{10:49:10 --- What the system believes, and how much it trusts it}

Fourteen seconds, and the densest phase in the run: the system refreshed what it
wants, reconciled the entities it knows about, ran its belief filter, and
updated how far it trusts each of its evidence channels.  Nothing external
happened.  This is the system forming an opinion about its own situation before
committing to anything --- and the patterns here are what make that opinion
inspectable rather than a mood.

\hypertarget{pat-r1}{}\paragraph{Belief State as Operational Hypotheses (R1).}\label{pat:r1} \textbf{If} you want decisions to depend on tractable hypotheses (progress, blockers, pattern fit, uncertainty) rather than raw conversation history. \textbf{However} if the belief state is just the conversation, the agent cannot cleanly update, compare, or falsify its own assumptions. \textbf{Then} represent the belief state as a map of operational hypotheses
\begin{quote}
  (\texttt{:goal}, \texttt{:subgoals}, \texttt{:blockers}, \texttt{:pattern-fit}, \texttt{:expected-next-observation}, \texttt{:uncertainty})
\end{quote}
updated each step using observation and outcomes. \textbf{Because} a compact belief state makes predict~$\rightarrow$~observe~$\rightarrow$~update implementable and testable.\worknote{IMPROVED: two admitted witness claims bind
the belief-state equation to production trace rows (a carried and a
bootstrapped coordinate row), each a Lean proposition checked at the
pinned record.}

\paragraph{Structured Observation Vector (R2).} \textbf{If} you want the agent's observation to be comparable across steps so scoring and learning do not collapse into prose interpretation. \textbf{However} tool output and user feedback arrive as heterogeneous artifacts (JSON, logs, diffs, snippets, natural language) that are hard to compare or weight consistently. \textbf{Then} maintain an explicit observation record as a typed, normalized feature map derived from session state (test status, compile status, diff size, failing-spec count, user-stated constraints, time since last anchor, contradiction flags). \textbf{Because} a stable observation vector is the prerequisite for meaningful precision control and policy scoring.\worknote{IMPROVED: the observe equation
carries an admitted witness at the empty-input reference model (fourteen
zeros --- the shape claim, proved not asserted).}

\paragraph{Predictive-Coding Belief Update (R3).} \textbf{If} you want belief updates to be principled, tunable (one learning rate; a separate trust setting per evidence channel), and operationally robust. \textbf{However} without an explicit rule the agent either slams its beliefs to each new observation (over-trusts) or never moves them (over-anchors), and an uncertainty estimate with no floor under it decays toward false certainty. \textbf{Then} move each belief \emph{toward} what was actually observed, but only as far as the learning rate and the trust placed in that channel allow; and track how uncertain the belief is by keeping a running average of the size of recent misses, measured against the freshly-updated belief, never letting that uncertainty fall below a floor.  It is the \emph{size} of a miss that counts, not its direction --- an overshoot and an undershoot are equally informative about how well the channel is understood --- and the floor is a standing admission that no channel is a perfect instrument. \textbf{Because} moving in proportion to trusted error is what predictive coding means, and the floor keeps uncertainty in a range where the ambiguity half of the action score still says something.  \emph{Invariant to preserve:} the belief must be able to move, and its uncertainty must not be able to reach zero.\worknote{IMPROVED: fourteen admitted prediction-error
claims, including the refusal branches this paragraph implies (broken
model refuses, malformed observation refuses, unobserved is omitted,
variance floored at zero and below).}

\paragraph{Evidence Precision Registry (R7).} \textbf{If} you want the agent to discount low-trust evidence without discarding it, and to become stricter when outcomes are volatile. \textbf{However} without an explicit precision layer, evidence weighting becomes implicit prompt rhetoric and cannot be audited. \textbf{Then} keep a per-channel \emph{precision registry} --- one trust weight per kind of evidence --- keyed to observation features, and apply it wherever prediction error, scoring, or gating occurs. \textbf{Because} making precision explicit turns trust into a tunable control surface rather than an undocumented bias.\worknote{IMPROVED: ten admitted precision
claims cover the still-error window family through the floor being
reached.}

\paragraph{Present-Fit Mismatch as a Consumed Quantity (R8).} \textbf{If} you want a number, comparable across time, that says how badly the agent's current belief explains what it just observed --- for validation, diagnostics, and alarms. \textbf{However} a mismatch number that nothing downstream consumes is not observability but decoration, and it decays into a false credential: this system computed exactly the per-tick scalar the first version of this entry recommended --- a trust-weighted average of how badly each channel missed --- and a source audit found it bound once, tagged for the trace, and read by nothing after scoring, so it was retired under a ruling (2026-09-01) rather than kept as a number to point at. \textbf{Then} compute the mismatch as a pure function of belief, observation, and precision, and emit it at the grain where a named consumer reads it.  Here that grain is the policy: the per-policy free energy $F(\pi)$, the fit of the observations under each candidate policy, admitted as a machine-checked claim at consecutive records of the pinned 2026-09-01 run. \textbf{Because} this mismatch is the quantity active-inference agents work to minimise, and it earns its place in the trace the same way any claim in this catalog does --- by something depending on it: a validation target that should fall as the agent learns, an alarm that fires when trusted evidence disagrees with the model, or a posterior that imports it.  \emph{Invariant to preserve:} every mismatch number in the trace has a consumer that can be named; a scalar with no reader is removed, not explained.\worknote{REWRITTEN (2026-09-14), the row-29
claim/model reconciliation: the per-tick scalar's producer was retired
2026-09-01 under a ruling (futon2 \texttt{5a66411};
\texttt{free\_energy.clj:7-12} records why), and the entry now teaches
what that audit found instead of recommending what was deleted. The
admitted evidence is registry row
\texttt{R8-policy-F-s4-consecutive-records-v1} (exact formal
proposition, s4 run 2026-09-01, receipts at
\texttt{runs/row-16-r8-policy-f-2026-09-12/}). Reopening a per-tick
scalar remains an operator decision, out of scope.}

\paragraph{Shared Kernel for the Predictive Forward Model (R4).} \textbf{If} you want the forward model's predictions to match what the live dynamics actually do, with one set of rules to maintain. \textbf{However} two implementations of the same dynamics is two places for bugs to hide and two surfaces where prediction can silently diverge from the live step. \textbf{Then} factor the live step into a \emph{pure deterministic kernel} (state, action, seed, opts)~$\rightarrow$~next-state plus a thin wrapper that threads defaults; the predictive forward model calls the same kernel and adds principled uncertainty. \textbf{Because} one implementation of the dynamics is the only place the rules live --- any update is reflected in both live and predicted behaviour.\worknote{IMPROVED: the forward model carries an
admitted witness at production pins (the float-carried v1.1 scope).}

\subsection*{10:49:26 --- Yesterday's debt outranks today's work}

Four seconds.  Before the system considered anything it might like to do, it
read its own outstanding obligations and found one: \texttt{repair-attempt-054},
a failure from an earlier run that an independent reviewer had rejected and that
had never been resolved.  That obligation pre-empted ordinary selection
entirely.  The system did not get to choose today's work until yesterday's was
discharged.

This is worth pausing on, because it is where the architecture's attitude to
failure is visible.  A failed attempt does not roll back and does not vanish.
It deposits a typed repair obligation \emph{into the field} --- into the same
landscape the loop reads when it decides what to do next --- so that a later
attempt encounters its own debt as ordinary terrain, and a stop-line rule makes
that debt jump the queue.  The loop never reverses; it only ever meets what it
previously broke.  What follows in this run is therefore not new work at all.
It is attempt-061 paying off attempt-054.

\subsection*{10:49:32 --- One hundred and eleven seconds of deliberation}

The longest phase in the run that did not involve waiting on another agent.  The
system proposed candidate actions, removed the ones the graph said were
infeasible, scored what remained, and recorded --- for each of the finalists ---
which term of the score governed the outcome and by how much.  It could also
have abstained here, and on many runs it does; on this one it did not.

The reason this phase takes nearly two minutes and produces a written
justification is the whole argument of the paper.  A system that merely picked
well would be a better scheduler.  A system that picks well \emph{and can be
asked why} is something a person can supervise.

\hypertarget{pat-r5}{}\paragraph{Expected Free Energy Inside an Auditable Controller (R5).}\label{pat:r5} \textbf{If} you want selection to balance canonical outcome risk and ambiguity with operational constraints. \textbf{However} calling every useful controller term ``EFE'' makes tuning indistinguishable from storytelling. \textbf{Then} compute and persist the canonical core separately --- \emph{risk}, how far the outcomes an action predicts sit from the outcomes the system prefers, plus \emph{ambiguity}, how uninformative those outcomes are expected to be --- and admit engineering terms only as a named augmentation, reported \emph{beside} that core and never folded into it. \textbf{Because} the decomposition makes clear whether a recommendation moved because of active-inference quantities or because of an engineering control.  \emph{Invariant to preserve:} the two-term core is persisted apart from everything wrapped around it, so any reader can ask which one moved the pick.\worknote{\wkflag{OUTSTANDING}: the ambiguity leg has no
general bridge between production's Gaussian entropy and the categorical
law --- proved: one isolated agreement at a single variance value, no
more. Decision on record: build a categorical estimator over the measured
Q and A, register the Gaussian path separately as what it actually
computes. Not yet built.}

\hypertarget{pat-r6}{}\paragraph{Candidate Pattern Action Space (R6).}\label{pat:r6} \textbf{If} you want pattern choice to be legible and comparable, and to avoid collapse into one habitual pattern. \textbf{However} if the action space is any pattern at any time, evaluation becomes unstructured and selection is dominated by recency. \textbf{Then} define the action set as a \emph{constructed} candidate set of pattern IDs produced by retrieval plus gating rules, of bounded size, with explicit reasons for inclusion and exclusion. \textbf{Because} constraining the action space makes downstream scoring meaningful rather than performative.\worknote{IMPROVED: policy-set and
policy-posterior (F-absent branch) witnesses admitted at a pinned
2026-09-12 production details record.}

\hypertarget{pat-r14}{}\paragraph{selection gain as Commitment Temperature (R14).}\label{pat:r14} \textbf{If} you want exploration versus exploitation to depend on operational pressure (deadlines, repeated failures, instability), not on randomness alone. \textbf{However} without an explicit dial the agent either thrashes or tunnels, and there is nothing to turn. \textbf{Then} maintain a \emph{commitment temperature} governing how decisively the best-scoring candidate wins: turned down, the system commits to its front-runner; turned up, it keeps the field open --- and couple that dial to diagnostic signals (recent regressions, uncertainty in the belief state, contradiction rate, time pressure). \textbf{Because} a single commitment dial allows systematic control of stochasticity without changing the scoring model at all.  \emph{Invariant to preserve:} exploration is a dial, not a random number generator buried in the selector.  \emph{Status in this implementation:} the dial is connected only under the declared full-score selection law. \srcref{futon2/src/futon2/aif/policy.clj} computes the temperature and the habit-adjusted score vector. When \texttt{:full-score-posterior} is requested and $F_\pi$ entered this tick, it selects the maximum full score and records \texttt{:habit-authority :live-in-selection-score}; otherwise it takes the head of R5's controller-score ordering and records \texttt{:habit-authority :counterfactual-only}. The decision record names the requested and applied laws and records a refusal when missing $F_\pi$ prevents the requested law from running. The dial is therefore live on one declared path and explicitly counterfactual on the fallback path, rather than silently changing meanings between them.\worknote{IMPROVED: a temperature witness is
admitted at a pinned trace row. The commitment-link constants are now
ruled (single factor one-half while any genuine trip is open, distinct
trip identities, no extra operator gate) with a sensitivity probe retained
over a 147-candidate production-shaped field. \wkflag{OUTSTANDING}: the composer
--- the ruled law is not yet wired into an applicable gain mode; a probe
showed the variational mode ignores task gain entirely, so wiring it there
would be fake evidence, and the integration must expose applicability.}

\paragraph{Hierarchical, Budget-Aware Action Selection (R11).\apparatusmark{}} \textbf{If} you want collective action selection that respects budget constraints at every hierarchy level --- never recommending allocations that exceed the aggregate per-tick budget, and factoring in downstream and coordination costs. \textbf{However} a naive union of independent selections leaves the budget invariant unenforced (over-subscription happens silently), while a single central coordinator throws away the value of local-agent factoring. \textbf{Then} resolve conflicts explicitly at \emph{each} hierarchy level: local agents propose within their sub-budgets, and a coordinator arbitrates the shared consumable budget across them. \textbf{Because} budget is a shared consumable, and enforcing the invariant at every level is the only way to keep multi-agent selection both locally factored and globally feasible. This is the \emph{resource} hierarchy --- how a shared per-tick budget is split across competing selections --- and it is orthogonal to the \emph{temporal / abstraction} hierarchy of R13 (a strategic loop choosing the task, a tactical loop rolling out the pattern-cascade that enacts it); the two compose but are not the same pattern.  \emph{Evidence:} the exact arbiter, mutually exclusive ranked-field adapter, over-subscription witnesses, and opt-in live selection boundary are tested together, including exact replay.  In Campaign S it selected one proposal from each live ranked field on both ticks, using all 7 units of the shared budget without exceeding it; exact replay reproduced both selections.\worknote{\wkflag{HONEST ABSENCE}: R11's lifecycle census records all
seven stages (commissioned through surfaced) typed absent, 0/7 --- item
16 of the row-24 negative-scope inventory. The 2026-09-13 refresh
attempt confirmed the census reads wm-trace records only and no
opportunity-counting procedure exists to extend; the recorded remedy is
build/integrate or rule non-load-bearing. Campaign S remains the
evidence; nothing newer exists yet.}

\paragraph{Hierarchical and Temporal Depth (R15).} \textbf{If} a fast tactical loop acts inside a slower strategic loop. \textbf{However} merely giving both levels a scalar score does not create a hierarchy. \textbf{Then} make the coupling explicit: strategic selection fixes the tactical target; witnessed tactical outcomes update the next strategic calibration state; and longer policies use a \emph{named} temporal discount of their own rather than overloading the commitment dial. \textbf{Because} hierarchy is a relationship between levels and timescales, not a second use of the same score.  \emph{Invariant to preserve:} each level has its own quantities, and no quantity does duty at two levels at once.  \emph{Evidence:} the opt-in two-tick campaign records both levels separately, rejects unwitnessed learning, folds a grounded fast outcome into a Beta update of the slow state, and uses that state to parameterise the next fast rollout.  Campaign S supplied the live positive branch: Tick A's independently approved and grounded outcome advanced the relevant slow posterior from Beta(1,1) to Beta(2,1), changed the slow mode from exploration to consolidation, and parameterised Tick B under that new state.  This is evidence for the explicit coupling, not for a general nested generative model.\worknote{IMPROVED: the machineAction
internal-branch behaviour that tactical selection rests on is now an
admitted supports-claim at captured machinery inputs (controller head,
first-max, last-max, abstain at the retained epsilon, tie control). Like
R11, the current period's R15 census rows are honest 0/7 absences; the
abstain-epsilon retention needed to change that is deployed.}

\subsection*{10:51:25 --- From a choice to a construction}

Two seconds.  Having chosen \emph{what} to do, the system assembled \emph{how}:
a bounded cascade of patterns fitted to the selected target.  A cascade is not a
to-do list.  It is a small construction --- some patterns feeding others, many
of them overlapping rather than sequential --- that can be checked before
anything runs, and that says which warrants are being combined and in what
order.

\hypertarget{pat-r13}{}\paragraph{Temporal Depth Beyond the Greedy Step (R13).}\label{pat:r13} \textbf{If} you want the agent to plan --- to score whole pattern cascades as policies, not single steps, and to prefer an action whose payoff is two or three ticks out over a greedy one-step win. \textbf{However} a rollout that looks one step ahead demonstrates the machinery but not depth; without a temporal hierarchy the horizon stays shallow. \textbf{Then} score the whole cascade over a horizon of more than one step, and introduce a two-timescale hierarchy: a slow (strategic) loop that updates priors and preferences, and a fast (tactical) loop that acts against them.  Scoring a route means adding up the costs of its steps with the later ones discounted, so that a distant payoff can matter without swamping the immediate situation. \textbf{Because} active inference plans over a \emph{policy} --- here, the pattern cascade for the selected mission --- not over the next action alone.  \emph{Invariant to preserve:} the thing being scored is the cascade, not its first move.\worknote{IMPROVED: a depth witness is
admitted at the 2026-09-12 machinery capture. \wkflag{OUTSTANDING}: production
resolves horizon nil in the wild, so the depth arm has never fired live;
the qualifying-run configuration (ruled) requires effective horizon at
least 2, which arms it.}

\paragraph{The Semilattice-Rollout Execution Witness (R16 realised, for cascades).\apparatusmark{}} \textbf{If} you want a mission's build phase (\textsc{instantiate}) to close only when its selected pattern cascade has actually been \emph{constructed} --- ``a live instance'', not a rich prose section describing one. \textbf{However} the honest-looking witness is a tautology in disguise: re-run the same fold that proposed the wiring and check it reproduced itself (a mirror, R16), or --- the failure we hit --- execute the cascade as a strict linear \emph{pipeline} through a fixed rule table, which silently abstains on every mission whose patterns that table does not already encode, so a genuinely well-posed cascade reads as unbuildable. \textbf{Then} take seriously that a cascade is a \emph{semilattice, not a chain} --- patterns compose both by sequence (one enables the next) and by overlap (one pattern belonging to several sub-constructions at once), and in practice the overlaps outnumber the sequences --- and execute it by an \emph{independent} rollout over that structure --- give each pattern a move interface whose \emph{consumes} are the salient tokens of its \textbf{If}/\textbf{However} clause and whose \emph{produces} are the tokens of its \textbf{Then}; admit a move only when it chains onto the frontier the earlier moves have built; and read discharge as the fraction of the mission's want-signature the assembled construction covers.\footnote{The want-coverage used here is a token-level \emph{proxy} for true discharge; training the move-prior on a structural-quality label (which cascades historically went well) is its natural strengthening, and turns the same rollout into a first estimate of the policy value R13 asks for. The executor is \srcref{futon3a/holes/labs/M-memes-arrows/rollout_execute.py} (over the \texttt{cascade\_rollout} move interface); the witness it earns is read by \srcref{futon3c/src/futon3c/logic/mission_clean.clj} (\texttt{instantiate-witness?}).} Mint the typed witness (\texttt{<mission>.executed.edn}) only when the construction both \emph{assembles} (at least one box was built) and \emph{discharges} (it covers some part of what the mission wanted), recording how much it covered as the realised signal; otherwise the phase re-opens. \textbf{Because} the rollout reconstructs the wiring from the patterns' \emph{own text} --- a different computation from the phylogeny-edge fold that predicted it --- so realised $\ne$ expected by construction: the witness is external to the predictor and cannot be farmed, which is precisely what R16 demands of a closed loop and what the strict-pipeline and re-run-the-fold alternatives both fail to give.

\subsection*{10:51:29 --- Handing the work to an agent that can fail}

Here the harness stops deliberating and becomes a client.  It dispatched the
work to a coding agent (\texttt{claude-7}) and waited five minutes and ten
seconds --- by far the longest single phase --- for a substantive commit.

Then the interesting part.  The build gate examined what came back and
\emph{flagged} it.  The system did not accept the result, and it did not throw
the attempt away either: it spent one bounded cure turn dispatching the same
author to fix its own build, which produced commit \texttt{6c537c6} eight
seconds later.  That single retry is the only backward motion permitted inside
an attempt.  It is bounded deliberately, because an unbounded repair loop is how
an agent harness quietly converts a failure into an expensive failure.

Note what the harness is not doing.  It is not writing the code, not reviewing
its own build, and not deciding that the commit is good.  Every load-bearing
quantity sits \emph{around} the coding agent rather than inside its prompt ---
which is the correction this architecture makes to its own predecessor, where
asking the agents to perform the formalism in-band mostly got in their way.

\subsection*{10:57:06 --- A different agent reviews}

The system dispatched a reviewer --- \texttt{codex-1}, a different agent, from a
different model family than the author --- and waited two minutes and eight
seconds for a verdict.  It came back \textsc{approve}, on first review.

The separation is not politeness.  An author asked to assess its own work
produces exactly the kind of evidence the whole apparatus exists to refuse, and
the same applies one level up: a reviewer's \emph{claim} to have reviewed is
itself self-manufactured, which is why review execution is corroborated from a
job-event stream the reviewer does not control before any review is allowed to
count.  That refinement was not designed in advance.  It was forced by two
measured failures in the first cohort, and R9 above carries the revision.

This is also where the difference between two kinds of calibration becomes
concrete.

\paragraph{Two-Layer Calibration (R12).} \textbf{If} predicted-versus-realised pairs are your evidence that the model is calibrated. \textbf{However} when the model's own output composes the field it is scored against, the comparison is self-referential: it tests dynamics-consistency, not whether the model's \emph{value} is right --- a model can pass forever while being worthless at predicting real outcomes. \textbf{Then} name and keep both layers explicitly: \emph{Layer~1} compares the model's predictions against its own later scores --- cheap, run every cycle, and permanently labelled never-value-evidence --- while \emph{Layer~2} compares predictions against witnessed outcomes the model did \emph{not} produce, which are scarce and therefore gated.  Only Layer~2 clears value, and Layer~1's cheap pass may never be reported as Layer~2 clearance. \textbf{Because} the only thing that escapes ``prior~$=$~value'' is a referent the model had no hand in; Layer~1's referent is its own output, so it can establish internal consistency and nothing more, while Layer~2's is external --- the one layer that can be \emph{wrong} about value, which is exactly why it is the gate.  \emph{Invariant to preserve:} the two layers are never allowed to share a label.  \emph{Evidence:} a deterministic report preserves the two records and their provenance under distinct fields and fails closed for a missing, failed, author-controlled, or self-witnessed Layer~2.  Campaign S produced two successful live records: on both ticks Layer~1 matched proposal to realised action, while Layer~2 was supplied by independent review and grounded re-observation; both combined gates cleared without substituting the cheap layer for the external one.

\subsection*{10:59:22 --- The witness, and the dial that moved}

Sixteen and a half seconds, and the only phase that changes the world outside
the model.  The approved commit was written into the external store as a typed
implementation-and-discharge record, and the store was then re-read through an
independent path to confirm the record was really there.  The trace line that
matters is \texttt{dial-moved?\ true}: a quantity the system does not author
read differently after the act than before it.

Everything preceding this phase is, strictly, a well-organised opinion.  This is
the phase that makes the loop a loop.

\hypertarget{pat-r16}{}\paragraph{Grounded Actuation, Not Re-Observation (R16).}\label{pat:r16} \textbf{If} you want a genuinely closed action--perception loop: the agent acts, and the world it observes on the next tick is materially different \emph{because} it acted. \textbf{However} this is seductively easy to fake --- an enactor that re-runs its own construction and measures how faithfully it reproduced its own output produces a perfectly honest number over a quantity that cannot move (a mirror, not an act). \textbf{Then} close the loop only when the act writes a \emph{typed, ungameable witness to a substrate outside the model} --- an endpoint proven inhabited by a query derived from the authored bindings, not from builder claims --- \emph{and} that witness feeds the next belief.  The witness demand must be \emph{shape-aware}: a deposit or a specification can be the typed, externally-checkable artifact when the act is data- or spec-shaped, because a binding that insists on an in-attempt code commit regardless of shape punishes honest authors who have nothing bindable (measured as \texttt{:artifact-binding-mismatch}, four occurrences in the first cohort; in the worst measured case the demand \emph{induced} a fabricated code change that independent review then rejected --- the guard against fabrication become a generator of it).  And \emph{typed} means typed over a population: every witness edge declares an enumerable domain and range, or loudly states that either cannot be enumerated, and a claim to cover the whole domain carries an executable coverage falsifier --- a four-entry map named \texttt{sampled-candidate-cleans} covers its declared sample 4/4 and claims nothing population-wide, where the same map under a whole-domain name passed shape checks while covering none of the 96 observed open-mission selections.  A witness contract without this population check can certify a fixture as the live substrate. \textbf{Because} the point of a perception--action loop is that action changes the world and the changed world changes the next belief; if ``realized'' is coverage-over-reproduction, the world never changes and the belief never updates.  \emph{Status in this implementation:} the enactor in this build is the mirror the pattern warns against. \srcref{futon2/src/futon2/aif/enact.clj}, lines 130--160 (at futon2 \texttt{e76c51c}), shells a fold engine and returns a typed in-model result --- ``\texttt{:constructed-wiring}, \texttt{:no-construction}, or \texttt{:enactment-failed}'' --- a construction \emph{inside} the model, not an effect outside it; two independently role-played endpoints reached that conclusion from opposite sides, R16 finding no outward action and R2 refusing an enactor's wiring as \emph{self}-report under the R9 verdict.  The pattern named this failure in advance, and the implementation is on the wrong side of it.\worknote{SHARPENED, both directions: the
admission this paragraph makes is now machine-checked --- machineAction
carries an admitted REFUTES claim at the live-selector scope (divergence
proved at a pinned production record) and, at a disjoint scope, an
admitted supports-claim for its internal branches. The mirror finding
stands. IMPROVED: live selection-enaction capture (entity state at close,
outcome entity, cross-ledger identity) is deployed; rows 14/23 wait on
fresh build-phase attempts to feed it.}

\subsection*{10:59:53 --- Closing at thirteen minutes and ten seconds}

The rest is bookkeeping, and it is short: the repair obligation from
attempt-054 was marked resolved (2.1s), that resolution was itself validated
against the substrate rather than taken on the loop's word (1.6s), an item was
queued for a human to adjudicate the next morning (4.5s), and the attempt closed
as \texttt{grounded-change}.

Thirteen minutes and ten seconds, three agent turns, one commit, one repaid debt, one item for a
person to look at.  That is a good run.  Most are not, and the honest accounting
of how many are not is in the companion paper --- but this is the shape the
system is trying to produce, and every pattern above exists because some earlier
version of this run went wrong in the way that pattern describes.\worknote{
The close-phase machinery this paragraph narrates has since been extended:
attempt close now retains the selected entity's state at close and the
outcome entity, and cross-ledger identity joins by run id only --- deployed
2026-09-12, awaiting fresh build-phase attempts to produce records (rows
14/23).}

\subsection*{Between runs: what the attempt left behind}

Some of the system's most consequential work cannot happen inside thirteen
minutes.  Attempt-061 deposited an adjudicated outcome, and that outcome joins
an accumulating record which is periodically re-read \emph{offline}, when
nothing is being decided --- to merge concepts that have turned out to be the
same, to prune structure that was never load-bearing, to retrain which
compositions get proposed in the first place, and occasionally to admit that no
existing pattern covers what keeps happening.

The patterns in this section are the slow ones.  They are also the ones where
embeddings, absent from the entire story so far, finally become essential.

\paragraph{Structure Learning by Bayesian Model Reduction (R17).} \textbf{If} you want the agent to form concepts and reorganise its generative model --- merge look-alike states, prune redundant connections, generalise across contexts --- from accumulated experience but no new data, offline.\footnote{The payoff is not merely fewer patterns; it is better search. If the system knows that several patterns form a useful constellation, it can propose richer candidate moves than a greedy one-pattern-at-a-time chooser.} \textbf{However} simply accumulating experience cannot do this: piling up counts only sharpens the mappings that already exist, and never removes a connection or ties two states together. \textbf{Then} run Bayesian Model Reduction over the accumulated counts --- ask whether the evidence \emph{already gathered} is better accounted for under a simpler model's assumptions, and accept the simplification only when the evidence favours it by roughly twenty to one.  No new data is collected: the same counts are re-read under the simpler assumption, which is what makes the operation cheap enough to run offline over the whole catalogue. \textbf{Because} model evidence --- not more data --- is what reorganises concepts, and this is the closed-form, biologically-plausible (synaptic pruning during sleep) operator that scores a reduced model against its parent.  \emph{Invariant to preserve:} merging requires a stated evidence threshold, and looking alike is never sufficient.  \emph{Evidence:} the offline envelope records the frozen parent, every proposed merge, the authoritative $\Delta F$ threshold and decision, the reduced posterior, and either a structural change or a principled non-change; its canonical input replays to exact equality.  Campaign S's live inter-tick phase froze 26 capability--mission observations, evaluated all 21 pairwise proposals among seven capability concepts in 20 milliseconds, replayed exactly, and carried its result into Tick B's proposal field.  Its accepted reductions collapsed all seven concepts into one, so the mechanism is run-evidenced while the usefulness of this particular reduction remains an operator-facing question rather than an automatic deployment claim.\worknote{IMPROVED: Dirichlet accumulation is
now an admitted witness (three-tick IEEE-residual correspondence at
production pins). HONEST STOPS, retained as typed refusals: the
capability-typed substrate is empty (no accumulation inputs live yet) and
the historical reduction's parent model was not retained --- both recorded
rather than reconstructed.}

\paragraph{Posterior Spread as an Exploration Control (R17\hspace{.1em}$^{\prime}$).} \textbf{If} you want unresolved structural hypotheses to remain eligible for exploration. \textbf{However} posterior standard deviation is uncertainty, not expected information gain: it contains no expectation over the observation and posterior update that a policy would induce. \textbf{Then} expose the spread under the honest name \emph{model-uncertainty bonus}, keep it outside the risk-plus-ambiguity core, and reserve the words ``information gain'' for a calculation that actually earns them. \textbf{Because} an honest exploration heuristic is useful, while a falsely canonical one makes the controller impossible to audit.  \emph{Invariant to preserve:} a quantity is named for what it computes, not for the role you wish it played.

\paragraph{Embedding-Nearness as the Structure-Learning Generator (R17\hspace{.1em}$^{\prime\prime}$).} \textbf{If} you want to run structure learning (R17) --- merging look-alike states, pruning redundant connections --- but need a tractable set of candidate reductions to test, because the space of possible merges is far too large to enumerate. \textbf{However} the control loop (R1--R16) is no help here: it never needed to know which patterns are alike, and hand-picking the candidates by fiat smuggles in the very structure you are trying to learn. \textbf{Then} exploit the fact that design patterns are written as text --- push the collection into a high-dimensional vector-space embedding and read off nearness, so that patterns sitting close become candidate merges; hand those to R17, which accepts only the mergers its evidence threshold clears, and let R17$^{\prime}$'s posterior spread rank which neighbourhoods to resolve first. \textbf{Because} the embedding is a cheap, model-free \emph{generator} of structural hypotheses while the evidence test is the \emph{disposer}: embeddings are inessential to the control loop but fundamental here, as its generator.  Embedding proposes; free energy disposes.  \emph{Invariant to preserve:} whatever proposes the merges must not also be what accepts them.

\hypertarget{pat-r17mint}{}\paragraph{Pattern Genesis from Evidence-Bearing Holes (R17\hspace{.1em}$^{\prime\prime\prime}$).}\label{pat:r17mint} \textbf{If} repeated, independently recorded episodes, failed retrievals, or fold policy-holes show that no existing pattern adequately expresses a recurring way of acting. \textbf{However} R6 can retrieve only the patterns already written down, R13 can compose only the moves already in the vocabulary, embedding nearness proposes similarity rather than a new action, and structure learning can reduce or compare the structures it has but cannot author a missing one. \textbf{Then} treat the gap as a proposal to \emph{enlarge the vocabulary}: bundle the triggering circumstances and policy-hole, supporting and challenging memories, rejected near-matches, and a falsifiable predicted outcome into a candidate IF/HOWEVER/THEN/BECAUSE pattern; mint it as a new move at an explicitly unearned prior; make concrete, through a fold, what enacting it is expected to change; and admit cascades containing it to witnessed enactment.  Only repeated R16-grounded outcomes and review can promote it; failures revise, merge, or prune it through the rest of R17. \textbf{Because} structure learning must sometimes enlarge the vocabulary in which policies can be expressed, not only rearrange or simplify an inherited vocabulary. Smith et al.'s active-inference model supplies the close theoretical precedent: insufficient state--outcome mappings motivate state-space expansion, after which reduction removes unnecessary structure \cite{smith2020structure}. Our control-space move is an explicit extension of that precedent, not a claim that their model authors design patterns.

\ifdefined\plopearlyoverview
\else
\emph{Worked instances.} The pattern \pid{structure/two-projections-of-one-quantity} was cited before it existed, caught as a policy-hole, minted, and then used to discharge the originating hole.  The pair \pid{process-coherence/status-refresh-before-work} and \pid{process-coherence/is-this-even-a-problem} distilled a measured failure family in which hundreds of expensive ticks worked from stale banners instead of reading the primary record.  The Zai memory pilot supplies the next, finer-grained test: several Lean episodes fit existing tactic patterns only loosely, so the Librarian must be able to return ``no adequate pattern'' and propose a narrower candidate rather than force every memory into its nearest neighbour.  Conversely, a volatile fact such as tool availability remains a reference memory rather than becoming a pattern.  These examples restate the enactment-and-review definition from the opening in operational terms: a hole warrants a candidate, not its promotion.\footnote{The older live route is recorded as \texttt{library-gap-findings $\rightarrow$ authoring $\rightarrow$ authored-seeds $\rightarrow$ pattern-library} in \srcref{futon2/holes/wm-pipeline-wiring.edn}; the generalized mint lane and the two-projections precedent are in \srcref{futon2/holes/labs/slush-demo/SPEC-full-loop-gfn.md}; typed-memory evidence and pattern-conditioned recall are developed in \srcref{futon3c/holes/missions/M-typed-memories.md}.}
\fi

\paragraph{Discharge-Trained Cascade Proposal (closing R16 $\rightarrow$ R6).} \textbf{If} you want the candidate cascades the agent proposes (R6) to improve with experience --- to reach for the compositions that have actually gone well, not a retrieval ranking frozen at authoring time. \textbf{However} embedding-similarity retrieval is a \emph{prior}, not a value: it ranks patterns that \emph{look} relevant to the problem text, and on the historical corpus the patterns a successful mission actually used sit hundreds of places down that cosine ranking --- similarity commits to nothing, so a proposal built from it alone cannot learn. \textbf{Then} train the proposal distribution from \emph{realised outcome}: score a completed mission's pattern-composition by the structural \emph{aliveness} of what R16 witnessed was really built --- Salingaros's measure of a structure's life, the product of its intensity and its internal coherence --- credit that aliveness back to the patterns the mission used, and fit a sampler that generates compositions in proportion to that reward --- a generative flow network over the cascade-construction semilattice, which samples the \emph{diversity} of alive compositions rather than collapsing to the single modal one a greedy prior returns. \textbf{Because} the outcome of using the catalogue is the only signal that separates relevance-that-commits from relevance-that-merely-retrieves; feeding the realised aliveness of built cascades back into how they are proposed is what makes the catalogue a \emph{policy that learns which of its own compositions to reach for}, not a static menu selected from by a fixed rule.\footnote{This pattern is the one with the most measured history behind it, and the history is mixed.  Credit trained on aliveness does separate the compositions that went well from those that did not, and survives a shuffled-label control.  But in a preregistered head-to-head over forty externally adjudicated builds, the trained proposer did \emph{not} beat the retrieval-prior incumbent on success rate, and a mission-conditioned variant made transfer to unseen missions worse rather than better.  What it did win on decisively was \emph{diversity} of proposals.  The disposition we adopted is therefore the honest one: keep it as a diversity supplier, not as a better chooser, until the supply of adjudicated labels is wide enough to settle the question.  The full protocol, metrics, and negative results are reported in the companion paper.}  \emph{Status in this implementation:} the running seam reads a trained artefact; it neither trains nor samples at live-loop time, so what is deployed is the diversity supplier, and the closed learning loop --- a witnessed outcome updating the proposal distribution and a later live selection consuming the update --- remains a proposal.\worknote{UNCHANGED: diversity-supplier
disposition holds; no new training or live sampling was added in the
2026-09 campaign.}

\paragraph{Capability Zones over the Embedded Landscape (staged; the self-model's second use of the metric).} \textbf{If} the agent must choose among its \emph{own capability gaps} --- the proposer surfaces one ``learn this action-class'' candidate per class it cannot yet enact, and the loop must rank them. \textbf{However} gap classes carry no outcome statistics \emph{by definition}: nothing about them has ever been enacted, so every honest posterior sits at its prior. This is measured, not hypothetical: after a night of four consecutive externally-witnessed groundings, the leading five candidates were all learn-action-class entries at bit-identical scores (fourteen such candidates, every per-class posterior still sitting untouched at its prior, zero recorded updates), and the non-discrimination tripwire refused the pick --- correctly, since preferring one never-observed option over another would be invented, not inferred. \textbf{Then} give the choice ground the agent already stands on: partition the same embedded landscape the stack renders as its mission field into \emph{action-class zones} --- a Voronoi tessellation over class seeds embedded in the same space --- and value each gap by what its zone \emph{carries}: structural load (enacted magnitude and summed turn-attestation of the features inside it), live demand (agent dwell-time in the zone, read from the running map), and posterior variance once months of repository history are described in capability terms and replayed into the per-class posteriors. \textbf{Because} capability in this stack is \emph{stigmergic} --- it lives in the built structure, not in agent biographies --- the structure's own load map is the only non-invented ground for preferences over learning; here the embedding stops being a generator of merge hypotheses and becomes the coordinate system the self-model's preferences run on. R17$^{\prime\prime}$'s scope claim extends here only as staged design: embeddings remain inessential to acting (R1--R16), and their role at this second learning-frontier use is prototyped, not established.\footnote{Chartered, not yet built: \srcref{futon0/holes/missions/M-capability-zones.md}, discharging the two stop-the-line obligations the measured refusal opened. Mechanism files: \srcref{futon2/src/futon2/aif/action_proposer.clj} (\texttt{gap-actions}), \srcref{futon2/src/futon2/aif/intrinsic_values.clj} (per-class Beta posteriors), \srcref{futon3c/src/futon3c/live_efe_map.clj} (the live coordinate field).} A dimensional note with a principled core: the \emph{operative} partition should live in a 3-D reduction of the embedding --- not the print-flat 2-D projection (whose adjacency distortions would silently misfile boundary features) and not the raw high-dimensional space (where distance concentration leaves every nearest-seed margin uniformly thin and the partition uninspectable). Three dimensions keep Voronoi cells compact and the tessellation \emph{walkable} --- and the deeper reason is an acceptance invariant in the spirit of R9: the partition the operator visually accepts must be the \emph{same object} the machine's preferences run on, so the machine must operate in a space a human can actually inspect, with the high-dimensional nearest-seed reading retained only as a boundary-distortion diagnostic.  \emph{Status:} the accepted 3-D projection disagrees with high-dimensional nearest-seed assignment on 79\% of commits, and no ablation establishes necessity; the coordinates are a staged prototype, and this paper makes no ground-truth or necessity claim for them.\worknote{\wkflag{STAGED}, and now formally excluded:
the 2026-09-13 paper exclusion list bars citing staged apparatus as
validated build capability; this entry and Native-Currency Discrimination
below remain proposals with their measured refusals intact.}

\paragraph{Native-Currency Discrimination (the R17$'$--R12 edge).} \textbf{If} the selector must rank capability gaps by earned evidence, while the non-discrimination gate rightly reads only the \emph{core} of the score --- risk and ambiguity, nothing else --- so that a bolt-on credit can never hide a flat estimator. \textbf{However} feeding the posteriors is not enough --- this is measured, not hypothetical: with the per-class credits fully fed and visibly distinct in the ranked field, the core stayed bit-identical and the gate refused again, because by the score-provenance discipline that credit enters only the engineering augmentation --- and each naive completion fails its own way: pushing structural uncertainty into the ambiguity leg \emph{alone} makes every unexplored class uniformly repulsive (anti-epistemic); fabricating an information-gain term violates the honesty rule of R17$'$; and re-pointing the gate at the augmented score guts the gate. \textbf{Then} make the core discriminate through \emph{both of its legs at once, in its native currency}: per-class zone load --- the structure's earned mass --- enters the \emph{preferences} the risk leg scores against, so risk carries what the anthill \emph{wants}; and each class's own predicted-outcome uncertainty replaces the class-flat ambiguity placeholder, so ambiguity carries what the machine genuinely \emph{does not know} about the outcome of learning that class. An unexplored-but-loaded zone then scores low risk and high ambiguity, and the selector's trade-off between them is exploration-versus-exploitation expressed \emph{inside} the canonical decomposition --- no new term, the gate untouched, and at the empty store the behaviour degrades exactly to today's honest refusal (uniform predictive variance, zero load shift). \textbf{Because} preferences are exactly where wanting lives in expected free energy and outcome ambiguity is exactly what the ambiguity leg measures; the R17$'$ quarantine is a \emph{boundary}, not a blanket --- it bars \emph{structure-posterior} spread from impersonating expected information gain, while the predictive variance of the scored action's own outcome is the ambiguity leg's proper content. That is the edge this pattern names: uncertainty about \emph{model structure} stays outside the core as an honestly-labelled bonus (R17$'$); uncertainty about \emph{an action's outcome} belongs inside it, as measurement (R12).\footnote{Chartered as M-capability-zones slice 2; the two measured refusals it discharges are attempt-036 (prior-flat posteriors) and attempt-041 (credit-fed but core-flat) --- the same tripwire correct both times, for different reasons.}


\subsection*{Underneath all of it: watching the machinery run}

One last pattern, which is not part of the story above because it is not part of
any run in particular.  Everything described so far is the loop observing the
\emph{world}.  Nothing in it observes the loop.  Every pathology the system has
suffered --- wedges, livelocks, mixed code images, silently lost wakes --- was
found afterwards, by archaeology on a run that had already failed.

\hypertarget{pat-r20}{}
\paragraph{Interoceptive Tripwires (R20).}\label{pat:r20} \textbf{If} you have hardened the loop's constructed core (R1--R16) and want the machinery itself to be observable while it runs live, because every pathology so far --- wedges, livelocks, mixed code images, silently lost wakes --- was discovered by post-hoc archaeology on a completed broken run. \textbf{However} the loop's observation channel is task-level only: it observes the world it acts \emph{on} (job outcomes, verdicts, witnesses, QA), while its own machinery is a hidden state with no likelihood mapping; and simulation tests of the known invariants pass by construction the moment the known bugs are fixed, so neither channel can notice the failure modes nobody has imagined yet. \textbf{Then} weave a supplementary harness of invariant tripwires around the live run at its phase boundaries: each wire is a near-infinite-precision prior over \emph{machine trajectories} (turn conservation, ledger closure, review-commit provenance binding, append-only repair stores and their status lattice, wedge and livelock detectors, wall-clock overrides, loaded-code/file-code coherence, and an environment-alphabet wire on which any unknown job state trips rather than defaulting); a trip is surprisal that belief updating cannot explain away, and the mandated epistemic action is \emph{freeze, record, park, summon} --- write the frozen state as a durable report, open a typed stop-line through the loop's own repair machinery, park an investigation join and summon the outer reviewer, on a fail-safe ladder that only ever degrades toward recording, never toward blocking the run. Calibrate before arming: the wires must \emph{retro-trip} on the known incident corpus (a wire that does not fire on the known past is a dud laser), and every closed incident thereafter must deposit either a new wire or an entry on an explicit blind-spot map. \textbf{Because} this is interoception, not vision: like any sensorium it admits no completeness theorem --- it gets identifiability checks, a measured coverage number, and a blind-spot map instead --- but an organism with only a hand (R16) and a preference model (C) learns of poison from the stomach ache, and each hotfix repairs one meal; an organism with interoceptive nerves notices at the moment of violation, with the state still warm, and every incident leaves it permanently more observant. The trip stream is also precision evidence about the machine itself (R7): a wire whose trips keep being blessed as model revisions demotes itself down the action ladder, and a genuine trip should lower the system's confidence in its own machinery until the finding is discharged --- the noticing feeding the feeling.\worknote{IMPROVED: a production catch on
2026-09-13 --- the loaded-code coherence wire (T10) fired four times at
run start on real serving-JVM drift after a 500-commit day, trips retained
as durable records; the drift was repaired by reload and the wire's
reading was correct. \wkflag{OUTSTANDING}: the trip-to-commitment-dial edge stays
chartered; its constants are ruled (row 18) and its commissioning
requirements are written, but no composer exists yet.}

\ifdefined\plopearlyoverview
\else
\emph{Evidence:} thirteen wires are catalogued and twelve enabled, specified from the fix-ledgers of the loop's own review history; they retro-trip 5/5 on the reconstructed incident corpus. The channel's first live catch was --- as the no-self-certification lore predicts --- about itself: fifteen shadow-mode trips from one wire whose composition baseline had been captured mid-load, an observation of emptiness mistaken for an empty observation; the repair (a missing observation admits the first real one, exactly once) is a likelihood-mapping correction, made at zero cost \emph{because} shadow mode. Armed operation has produced zero false alarms; one wire (author--artifact binding, chartered from a live incident in which the pipeline reviewed a sha \emph{narrated} rather than \emph{measured}) guarantees that every commit reviewed is the commit authored. The two mis-precisioned firings on record --- a wedge detector reading truly but at the wrong precision, and a summon-path roster failure that degraded one rung down the fail-safe ladder with nothing lost --- were repaired through the contract's own \emph{revise-the-model} route, on the ledger, with independent-review evidence, like any other repair. The link from a trip to the commitment dial of R14 --- a wire firing should make the system less sure of itself --- is chartered and not implemented.
\fi

% FIGURE TODO (aif-wiring-brain.svg): add node R20 on the INSULA (the
% interoception seat in predictive processing) --- the wire mesh watching the
% loop's own machinery. Edges: R20 -> R7 (trips are precision evidence about
% internal channels); dashed chartered edge R20 -> R14 (interoceptive gamma).
% Styling: green role-outline (wired, armed, live); AMBER faithfulness dot
% (principled approximation: invariant-violation surprisal as interoceptive
% prediction error, one repair --- learned per-wire precision --- from
% canonical). Caption gains one sentence when the node lands.

\ifdefined\plopearlyoverview
\else
\bigskip

Figure~\ref{fig:aif-brain} places the whole catalogue on one diagram.  Reading
it alongside the run above: the core inference nodes sit on the brain regions
their active-inference roles are associated with, and the apparatus sits below
them.  The clockwise cycle through the core is the 10:49:10 and 10:49:32
episodes --- perceive, believe, evaluate, select --- and the descent into the
apparatus is everything from 10:51:29 onward, where the loop stops thinking and
starts having to prove things.  Two kinds of colour are doing two different
jobs, and the caption distinguishes them: the node outline reports how far
\emph{built} a component is, while the small dot on its top edge reports how far
\emph{faithful} the quantity is to the theory.  A node can honestly be fully
built and only analogically faithful; those are different admissions, and
collapsing them into one scale would hide the more important one.

\hypertarget{pat-r19}{}
\begin{figure}[h]
\centering
\includegraphics[width=\linewidth]{aif-wiring-brain}
\caption{The active-inference loop, drawn on a brain. Core inference nodes sit on their Friston brain regions; the implementation apparatus sits below. Node styling reports implementation status: green is wired, amber is partial, and blue is live on a sibling surface. Lavender is deliberately not another maturity grade: R19 marks the preference organ or ``belly'' --- what the system \emph{wants}, against which risk is scored --- built and steering through a provisional in-memory join while its durable proof join remains open. The small top-edge dots report a different property: how faithful each named quantity is to the theory. Green is derived from the free-energy principle; amber is a principled approximation one repair away from the canonical quantity; pink is analogical; and gray is an engineering quantity that makes no theoretical claim at all.  Reading the figure as a whole: the perception--action cycle runs clockwise through the core nodes, the apparatus below is what makes each step auditable, and the gray dots mark every place where we are steering with something the theory did not give us.}
\label{fig:aif-brain}
\end{figure}
\fi

%  Thus R19 can have a lavender role-outline and a green faithfulness dot without contradiction. At R5 the implemented Expected Free Energy ($\mathrm{risk}+\mathrm{ambiguity}$) is drawn distinctly from the multi-objective controller that wraps it. Active gray-dot controls retract onto AIF-theory roles with an explicit residual; telemetry-only gray quantities instead have a diagnostic referent but no load-bearing role in live selection (Table~\ref{tab:aif-retract}). R3's green styling is model-relative: production executes the exact categorical filter under its declared A/B/D model; empirical calibration of the hand-set A entries remains open.
\FloatBarrier
